\chapter{Υλοποίηση και Στήσιμο}
\label{ch:implementation}


Το παρόν κεφάλαιο περιγράφει αναλυτικά τη διαδικασία στησίματος του
\textlatin{Data Space} σε νέο μηχάνημα, από τη δημιουργία των
κρυπτογραφικών πιστοποιητικών έως τη φόρτωση δεδομένων και τη δημιουργία
πολιτικών εξουσιοδότησης. Τα βήματα που παρουσιάζονται πρέπει να
εκτελεστούν με την ακριβή σειρά που αναφέρεται, καθώς κάθε βήμα αποτελεί
προϋπόθεση για το επόμενο.

Όλο το υλικό και ο κώδικας της υλοποιήσης είναι διαθέσιμος στο αποθετήριο
\textlatin{GitHub} \href{https://github.com/sakakos/Data-Spaces}{\textbf{\textlatin{Data-Spaces}}}.

Για την εκτέλεση της υλοποίησης απαιτούνται τα ακόλουθα εργαλεία στο
μηχάνημα ανάπτυξης:

\begin{itemize}
    \item \textlatin{Docker Desktop} εγκατεστημένο και σε λειτουργία.
    \item \textlatin{Python 3} εγκατεστημένο τοπικά.
    \item \textlatin{Git} για κλωνοποίηση του αποθετηρίου.
\end{itemize}

% ------------------------------------------------------------------
\section{Βήμα 1: Δημιουργία Κρυπτογραφικών Πιστοποιητικών}
\label{sec:step1-certs}
% ------------------------------------------------------------------

Το πρώτο βήμα είναι η δημιουργία της κρυπτογραφικής υποδομής
(\textlatin{PKI}) που θα χρησιμοποιηθεί για την ταυτοποίηση των
συμμετεχόντων. Απαιτούνται τρία ζεύγη κλειδιών και πιστοποιητικών:
ένα για την τοπική Αρχή Πιστοποίησης (\textlatin{CA}), ένα για τον
πάροχο και ένα για τον καταναλωτή.

Το \textlatin{script} \textlatin{01\_generate\_certs.py} αυτοματοποιεί
την παραγωγή όλων των πιστοποιητικών χρησιμοποιώντας τη βιβλιοθήκη
\textlatin{Python} \textlatin{cryptography}:
\vspace{0.2cm}

\selectlanguage{english}
\begin{lstlisting}[language=bash, caption={Certificate generation script execution.}, stringstyle=\color{black}, keywordstyle=\color{black}]
cd "Setup Instructions"
python 01_generate_certs.py
\end{lstlisting}
\selectlanguage{greek}

Το \textlatin{script} παράγει τα ακόλουθα αρχεία:

\begin{itemize}
    \item \textlatin{ca.key}, \textlatin{ca.pem}: Το ιδιωτικό κλειδί και
    το αυτο-υπογεγραμμένο πιστοποιητικό της \textlatin{PlegmaCA}.
    \item \textlatin{client.key}: Το ιδιωτικό κλειδί του παρόχου σε
    μορφή \textlatin{PKCS\#8} (\textlatin{BEGIN PRIVATE KEY}).
    \item \textlatin{client\_rsa.key}: Το ίδιο κλειδί σε μορφή
    \textlatin{PKCS\#1} (\textlatin{BEGIN RSA PRIVATE KEY}). Η ύπαρξη
    δύο μορφών είναι αναγκαία γιατί το \textlatin{Kong} χρησιμοποιεί
    τη βιβλιοθήκη \textlatin{lua-resty-jwt} που απαιτεί
    \textlatin{PKCS\#1}, ενώ ο \textlatin{Keyrock} χρησιμοποιεί
    τη \textlatin{node-forge} που επίσης απαιτεί \textlatin{PKCS\#1}.
    \item \textlatin{client.pem}: Το πιστοποιητικό του παρόχου.
    \item \textlatin{fullchain.pem}: Η πλήρης αλυσίδα πιστοποιητικών
    (πιστοποιητικό παρόχου + πιστοποιητικό \textlatin{CA}).
    \item \textlatin{consumer.key}, \textlatin{consumer.pem},
    \textlatin{consumer\_fullchain.pem}: Τα αντίστοιχα αρχεία για
    τον καταναλωτή.
\end{itemize}

Δύο παράμετροι είναι κρίσιμες για τη σωστή λειτουργία:

\begin{enumerate}
    \item Η επέκταση \textlatin{keyUsage: digitalSignature} ορίζεται
    ρητά στα πιστοποιητικά παρόχου και καταναλωτή. Ο \textlatin{Keyrock}
    ελέγχει αυτή την επέκταση κατά την επαλήθευση πιστοποιητικών
    μέσω του μηχανισμού \textlatin{extparticipant}. Χωρίς αυτήν,
    η ταυτοποίηση αποτυγχάνει.

    \item Το αναγνωριστικό \textlatin{EORI} τοποθετείται στο πεδίο
    \textlatin{Serial Number} του \textlatin{Distinguished Name}
    κάθε πιστοποιητικού, ώστε ο \textlatin{Satellite} να μπορεί
    να αντιστοιχίσει ένα πιστοποιητικό στον αντίστοιχο συμμετέχοντα
    του μητρώου.
\end{enumerate}

Μετά την εκτέλεση, το \textlatin{script} εκτυπώνει τα αποτυπώματα
(\textlatin{fingerprints}) και τη συμβολοσειρά \textlatin{x5c} που
πρέπει να αντιγραφούν στα αρχεία παραμετροποίησης.

Τα παραγόμενα αρχεία πρέπει να αντιγραφούν στους κατάλληλους φακέλους:

\selectlanguage{english}
\begin{lstlisting}[language=bash, caption={Copying certificates to component directories.}]
# Copy to Kong/iShare/
copy *.key Kong\iShare\
copy *.pem Kong\iShare\

# Copy PKCS#1 key to Keyrock (node-forge requires this format)
copy client_rsa.key Keyrock\iShare\client.key
copy fullchain.pem  Keyrock\iShare\fullchain.pem
\end{lstlisting}
\selectlanguage{greek}

Η τελευταία εντολή είναι κρίσιμη: στον φάκελο \textlatin{Keyrock/iShare/}
αντιγράφεται το \textlatin{client\_rsa.key} (μορφή \textlatin{PKCS\#1})
ως \textlatin{client.key}, γιατί η βιβλιοθήκη \textlatin{node-forge}
του \textlatin{Keyrock} δεν αναγνωρίζει τη μορφή \textlatin{PKCS\#8}.

\subsection{Αντιμετώπιση \textlatin{Windows Line Endings}}
\label{subsec:line-endings}

Σε περιβάλλον \textlatin{Windows}, τα αρχεία \textlatin{PEM} μπορεί
να αποθηκευτούν με \textlatin{CRLF} (\textlatin{\textbackslash r\textbackslash n})
αντί \textlatin{LF} (\textlatin{\textbackslash n}) και να περιέχουν
\textlatin{BOM} (\textlatin{Byte Order Mark}) στην αρχή. Και τα
δύο προκαλούν σφάλμα \textlatin{«Invalid PEM formatted message»}
στον \textlatin{Keyrock}. Η λύση είναι η μετατροπή των αρχείων σε
\textlatin{UTF-8} χωρίς \textlatin{BOM} και με \textlatin{Unix line
endings}. Τα \textlatin{scripts} παραγωγής \textlatin{JWT}
(\textlatin{gen\_token.py} και \textlatin{gen\_consumer\_token.py})
χρησιμοποιούν ρητά
\textlatin{.replace('\textbackslash r\textbackslash n',
'\textbackslash n')} πριν την ανάγνωση πιστοποιητικών, ώστε
να αντιμετωπίζεται αυτόματα το πρόβλημα αυτό.

% ------------------------------------------------------------------
\section{Βήμα 2: Παραμετροποίηση \textlatin{iSHARE Satellite}}
\label{sec:step2-satellite}
% ------------------------------------------------------------------

Ο \textlatin{iSHARE Satellite} παραμετροποιείται μέσω του αρχείου
\textlatin{satellite.yml}. Στο αρχείο αυτό ορίζονται τρία κρίσιμα
στοιχεία: τα κρυπτογραφικά κλειδιά του \textlatin{Satellite} (τα
οποία ταυτίζονται με αυτά του παρόχου, καθώς λειτουργεί υπό την
ταυτότητα \textlatin{EU.EORI.NLPLEGMA}), η λίστα εμπιστευμένων
\textlatin{CA} (\textlatin{trusted\_list}) με το αποτύπωμα
\textlatin{SHA-256} της \textlatin{PlegmaCA}, και το μητρώο
εγκεκριμένων συμμετεχόντων (\textlatin{parties}) με εγγραφές για
τον πάροχο και τον καταναλωτή.

Κάθε εγγραφή στο μητρώο περιέχει το αναγνωριστικό \textlatin{EORI},
την κατάσταση (\textlatin{Active}), τους ρόλους και το αποτύπωμα
\textlatin{SHA-256} του πιστοποιητικού:

\selectlanguage{english}
\begin{lstlisting}[language={}, caption={Excerpt of satellite.yml for participant registration.}]
parties:
  - id: EU.EORI.NLPLEGMA
    name: Plegma Data Provider
    status: Active
    certifications:
      - role: IdentityProvider
      - role: AuthorisationRegistry
    certificates:
      - certificate_fingerprint: "19A0DD7A..."

  - id: EU.EORI.NLPLEGMACONSUMER
    name: Plegma Data Consumer
    status: Active
    certifications:
      - role: ServiceConsumer
    certificates:
      - certificate_fingerprint: "B26DB03F..."
\end{lstlisting}
\selectlanguage{greek}

Τα αποτυπώματα πρέπει να αντιστοιχούν ακριβώς στα πιστοποιητικά που
παρήχθησαν στο Βήμα 1. Αν αναγεννηθούν τα πιστοποιητικά, πρέπει
υποχρεωτικά να ενημερωθούν και τα αποτυπώματα στο
\textlatin{satellite.yml}.

% ------------------------------------------------------------------
\section{Βήμα 3: Παραμετροποίηση \textlatin{Kong API Gateway}}
\label{sec:step3-kong}
% ------------------------------------------------------------------

Το \textlatin{Kong} παραμετροποιείται μέσω του αρχείου
\textlatin{kong.yml} σε δηλωτική μορφή (\textlatin{DB-less mode}).
Ορίζεται μια υπηρεσία που δρομολογεί τις αιτήσεις προς τον
\textlatin{Orion-LD}, και προσαρτάται το \textlatin{plugin}
\textlatin{ngsi-ishare-policies} με τρεις ενότητες παραμέτρων:

\selectlanguage{english}
\begin{lstlisting}[language={}, caption={Excerpt of kong.yml -- plugin configuration.}]
plugins:
  - name: ngsi-ishare-policies
    config:
      jws:
        identifier: EU.EORI.NLPLEGMA
        root_ca_file: /iShare/certificate.pem
        private_key: /iShare/key.pem
        x5c: "<cert_base64>,<CA_base64>"
      satellite:
        identifier: EU.EORI.NLPLEGMA
        host: http://ishare-satellite:8080
        token_endpoint: /token
        trusted_list_endpoint: /trusted_list
      ar:
        identifier: EU.EORI.NLPLEGMA
        host: http://fiware-keyrock:3007
        token_endpoint: /oauth2/token
        delegation_endpoint: /ar/delegation
\end{lstlisting}
\selectlanguage{greek}

Η ενότητα \textlatin{jws} περιέχει τα κρυπτογραφικά στοιχεία του
παρόχου που χρειάζεται το \textlatin{Kong} για να υπογράφει τα
δικά του \textlatin{JWT} κατά την επικοινωνία με τον
\textlatin{Satellite}. Η ενότητα \textlatin{satellite} ορίζει πού
βρίσκεται ο \textlatin{Satellite} ώστε το \textlatin{Kong} να
μπορεί να επαληθεύσει αλυσίδες εμπιστοσύνης. Η ενότητα \textlatin{ar}
ορίζει πού βρίσκεται το Μητρώο Εξουσιοδότησης (\textlatin{Keyrock})
ώστε το \textlatin{Kong} να μπορεί να αναζητά πολιτικές εκχώρησης.

Η τιμή \textlatin{x5c} πρέπει να αντιστοιχεί στα πιστοποιητικά
παρόχου και \textlatin{CA} σε μορφή \textlatin{Base64}, χωρίς
τις γραμμές \textlatin{BEGIN/END CERTIFICATE}.

% ------------------------------------------------------------------
\section{Βήμα 4: Εκκίνηση \textlatin{Containers}}
\label{sec:step4-startup}
% ------------------------------------------------------------------

Η σειρά εκκίνησης είναι κρίσιμη. Το \textlatin{Kong}
\textlatin{docker-compose.yml} δηλώνει τα δίκτυα \textlatin{fiware\_default}
και \textlatin{fiware-keyrock\_default} ως \textlatin{external}, δηλαδή
απαιτεί να υπάρχουν ήδη πριν την εκκίνηση.

Η σωστή σειρά εκκίνησης είναι:

\selectlanguage{english}
\begin{lstlisting}[language=bash, caption={Docker network creation and container startup order.}, stringstyle=\color{black}, keywordstyle=\color{black}]
# 1. Create Docker networks
docker network create fiware_default
docker network create fiware-keyrock_default

# 2. Start Orion-LD (uses two compose files)
cd Orion-LD
docker-compose \
  -f docker-compose/common.yml \
  -f docker-compose/orion-ld.yml \
  --env-file .env \
  up -d mongo-db orion

# 3. Start Keyrock (wait 30s for MySQL to initialize)
cd ../Keyrock
docker-compose up -d

# 4. Start Kong (both networks must exist)
cd ../Kong
docker-compose up -d
\end{lstlisting}
\selectlanguage{greek}

Μετά την εκκίνηση, η εντολή \textlatin{docker ps} πρέπει να
εμφανίζει 7 \textlatin{containers}: \textlatin{fiware-orion-ld},
\textlatin{db-mongo}, \textlatin{fiware-keyrock}, \textlatin{db-mysql},
\textlatin{kong-standalone}, \textlatin{dsba-pdp} και
\textlatin{ishare-satellite}.

% ------------------------------------------------------------------
\section{Βήμα 5: Δημιουργία Εφαρμογής \textlatin{OAuth} στον \textlatin{Keyrock}}
\label{sec:step5-app}
% ------------------------------------------------------------------

Ο \textlatin{Keyrock} χρειάζεται μια εφαρμογή \textlatin{OAuth 2.0}
για να εκδίδει \textlatin{access tokens}. Σε νέα εγκατάσταση η
εφαρμογή δεν υπάρχει και πρέπει να δημιουργηθεί μέσω του γραφικού
περιβάλλοντος:

\begin{enumerate}
    \item Άνοιγμα \textlatin{browser} στο \textlatin{http://localhost:3007}.
    \item Σύνδεση με τα διαπιστευτήρια \textlatin{admin@test.com} /
    \textlatin{1234}.
    \item Δημιουργία νέας εφαρμογής (\textlatin{Applications} $\rightarrow$
    \textlatin{Register}) με τη ροή \textlatin{Client Credentials}
    ενεργοποιημένη.
    \item Σημείωση του \textlatin{Client ID} και \textlatin{Client
    Secret} που εκδίδονται.
    \item Ενημέρωση του αρχείου \textlatin{02\_setup\_keyrock.py} με
    τις νέες τιμές \textlatin{CLIENT\_ID} και \textlatin{CLIENT\_SECRET}.
\end{enumerate}

% ------------------------------------------------------------------
\section{Βήμα 6: Εφαρμογή \textlatin{Patches} στον \textlatin{Keyrock}}
\label{sec:step6-patches}
% ------------------------------------------------------------------

Ο \textlatin{Keyrock} απαιτεί δύο τροποποιήσεις στον πηγαίο κώδικά
του (\textlatin{patches}) για να λειτουργήσει σωστά ως Μητρώο
Εξουσιοδότησης \textlatin{iSHARE}.

\subsection{\textlatin{Patch 1}: \textlatin{model\_oauth\_server.js}}
\label{subsec:patch1}

Το πρώτο \textlatin{patch} προσθέτει το πεδίο \textlatin{admin} στη
λίστα \textlatin{attributes} που ανακτώνται από τη βάση δεδομένων
κατά την ταυτοποίηση χρήστη. Χωρίς αυτό, ο \textlatin{Keyrock} δεν
αναγνωρίζει τον διαχειριστή ως τέτοιον, και η δημιουργία πολιτικών
εκχώρησης μέσω \textlatin{POST /ar/policy} αποτυγχάνει με σφάλμα
\textlatin{403 Forbidden}.

Η εφαρμογή γίνεται μέσω \textlatin{docker cp} και
\textlatin{docker exec}:

\selectlanguage{english}
\begin{lstlisting}[language=bash, caption={Applying patch1 to Keyrock.}, stringstyle=\color{black}, keywordstyle=\color{black}]
docker cp Keyrock/patch1.js fiware-keyrock:/tmp/patch1.js
docker exec fiware-keyrock node /tmp/patch1.js
\end{lstlisting}
\selectlanguage{greek}

\subsection{\textlatin{Patch 2}: \textlatin{configService.js}}
\label{subsec:patch2}

Το δεύτερο \textlatin{patch} τροποποιεί τον τρόπο φόρτωσης
κρυπτογραφικών κλειδιών. Στην αρχική έκδοση, οι μεταβλητές
περιβάλλοντος \textlatin{IDM\_PR\_CLIENT\_KEY} και
\textlatin{IDM\_PR\_CLIENT\_CRT} αντιμετωπίζονται ως το ίδιο
το περιεχόμενο του κλειδιού. Το \textlatin{patch} τροποποιεί
αυτή τη λογική ώστε, αν η τιμή ξεκινά με \textlatin{/} (δηλαδή
είναι διαδρομή αρχείου), ο \textlatin{Keyrock} να διαβάζει
το αρχείο αντί να χρησιμοποιεί απευθείας την τιμή:

\selectlanguage{english}
\begin{lstlisting}[language=bash, caption={Applying patch2 to Keyrock.}]
docker cp Keyrock/patch_config.js fiware-keyrock:/tmp/patch_config.js
docker exec fiware-keyrock node /tmp/patch_config.js
docker restart fiware-keyrock
\end{lstlisting}
\selectlanguage{greek}

Μετά την εφαρμογή και των δύο \textlatin{patches}, ο
\textlatin{Keyrock} πρέπει να επανεκκινηθεί ώστε να ισχύσουν
οι αλλαγές.

Αξίζει να σημειωθεί ότι τα \textlatin{patches} αυτά πρέπει να
εφαρμόζονται κάθε φορά που ο \textlatin{Keyrock container}
δημιουργείται εκ νέου, καθώς οι αλλαγές γίνονται εντός του
\textlatin{container} και δεν αποθηκεύονται στο \textlatin{Docker
volume}. Σε περιβάλλον παραγωγής, οι τροποποιήσεις θα ενσωματώνονταν
σε μια προσαρμοσμένη εικόνα \textlatin{Docker} μέσω
\textlatin{Dockerfile}.

% ------------------------------------------------------------------
\section{Βήμα 7: Αρχική Εγγραφή Καταναλωτή}
\label{sec:step7-consumer}
% ------------------------------------------------------------------

Πριν από τη δημιουργία πολιτικής εκχώρησης, ο καταναλωτής πρέπει
να έχει πραγματοποιήσει τουλάχιστον ένα αίτημα προς τον
\textlatin{Keyrock}. Αυτό είναι απαραίτητο λόγω του μηχανισμού
\textlatin{extparticipant}: ο \textlatin{Keyrock} δεν δημιουργεί
εγγραφή χρήστη για εξωτερικούς \textlatin{iSHARE} συμμετέχοντες
μέχρι την πρώτη τους επικοινωνία. Χωρίς εγγραφή δεν υπάρχει
εσωτερικό \textlatin{UUID}, και χωρίς \textlatin{UUID} η πολιτική
εκχώρησης δεν μπορεί να αποθηκευτεί.

Η αρχική εγγραφή γίνεται στέλνοντας ένα \textlatin{iSHARE JWT} του
καταναλωτή στο \textlatin{endpoint} \textlatin{POST /oauth2/token}
του \textlatin{Keyrock}:

\selectlanguage{english}
\begin{lstlisting}[language=bash, caption={Initial consumer registration in Keyrock.}, stringstyle=\color{black}, keywordstyle=\color{black}]
CONSUMER_JWT=$(docker run --rm \
  -v "./Kong/iShare:/iShare" python:3.9-slim sh -c \
  "pip install PyJWT cryptography -q && \
   python3 /iShare/gen_consumer_token.py")

curl -s -X POST "http://localhost:3007/oauth2/token" \
  -H "Content-Type: application/x-www-form-urlencoded" \
  -d "grant_type=client_credentials \
      &scope=iSHARE \
      &client_assertion_type=...jwt-bearer \
      &client_assertion=${CONSUMER_JWT} \
      &client_id=EU.EORI.NLPLEGMACONSUMER"
\end{lstlisting}
\selectlanguage{greek}

Κατά την εκτέλεση, ο \textlatin{Keyrock} αποκωδικοποιεί το
\textlatin{JWT}, επικοινωνεί με τον \textlatin{Satellite} για να
επαληθεύσει τον καταναλωτή, δημιουργεί αυτόματα εγγραφή χρήστη στη
\textlatin{MySQL} με ένα νέο εσωτερικό \textlatin{UUID} και
επιστρέφει \textlatin{access token}.

Αν η απόκριση είναι σελίδα \textlatin{HTML} με σφάλμα αντί για
\textlatin{JSON}, οι πιο συνηθισμένες αιτίες είναι: τα
\textlatin{patches} δεν εφαρμόστηκαν, τα πιστοποιητικά έχουν
\textlatin{BOM} ή \textlatin{CRLF}, ή ο \textlatin{Keyrock}
δεν μπορεί να επικοινωνήσει με τον \textlatin{Satellite}.

% ------------------------------------------------------------------
\section{Βήμα 8: Δημιουργία Πολιτικής Εξουσιοδότησης}
\label{sec:step8-policy}
% ------------------------------------------------------------------

Μετά την αρχική εγγραφή, ανακτάται το εσωτερικό \textlatin{UUID}
του καταναλωτή από τη βάση δεδομένων \textlatin{MySQL}:

\selectlanguage{english}
\begin{lstlisting}[language=bash, caption={Retrieving consumer UUID from MySQL.}, stringstyle=\color{black}, keywordstyle=\color{black}]
docker exec db-mysql mysql -u root -psecret idm \
  -e "SELECT id FROM user \
      WHERE username='EU.EORI.NLPLEGMACONSUMER';"
\end{lstlisting}
\selectlanguage{greek}

Κρίσιμη λεπτομέρεια: η πολιτική αποθηκεύεται με αναφορά στο
εσωτερικό \textlatin{UUID} του καταναλωτή (π.χ.
\textlatin{44d5c69f-0e2e-4565-a661-9c2e74f4ac1f}) και όχι στο
\textlatin{EORI} του. Αυτό το \textlatin{UUID} αλλάζει σε κάθε
νέα εγκατάσταση, οπότε η πολιτική πρέπει να ενημερωθεί αντίστοιχα.

Η δημιουργία της πολιτικής γίνεται μέσω \textlatin{POST /ar/policy}
με \textlatin{token} διαχειριστή:

\selectlanguage{english}
\begin{lstlisting}[language=bash, caption={Creating the delegation policy.}, stringstyle=\color{black}, keywordstyle=\color{black}]
# Get admin token
ADMIN_TOKEN=$(curl -s -X POST \
  "http://localhost:3007/oauth2/token" \
  -d "grant_type=password \
      &username=admin@test.com&password=1234 \
      &client_id=<CLIENT_ID> \
      &client_secret=<CLIENT_SECRET>" \
  | jq -r .access_token)

# Send policy from file
curl -s -X POST "http://localhost:3007/ar/policy" \
  -H "Content-Type: application/json" \
  -H "Authorization: Bearer ${ADMIN_TOKEN}" \
  --data-binary "@Keyrock/policy.json"
\end{lstlisting}
\selectlanguage{greek}

Η πολιτική που δημιουργείται επιτρέπει στον καταναλωτή να εκτελεί
λειτουργίες \textlatin{GET} σε οντότητες τύπου
\textlatin{HouseholdElectricMeasurement}. Η δομή της ακολουθεί
το μοντέλο \textlatin{delegationEvidence} του \textlatin{iSHARE}
και περιλαμβάνει τον εκδότη (\textlatin{policyIssuer}), τον αποδέκτη
(\textlatin{accessSubject}), τους επιτρεπόμενους πόρους και ενέργειες,
και το χρονικό παράθυρο ισχύος.

Σε περιβάλλον \textlatin{Windows}, το σώμα \textlatin{JSON} πρέπει
να σταλεί μέσω αρχείου (\textlatin{--data-binary "@αρχείο"}) και όχι
ως παράμετρος γραμμής εντολών, για αποφυγή προβλημάτων
κωδικοποίησης \textlatin{BOM}.

% ------------------------------------------------------------------
\section{Βήμα 9: Φόρτωση Δεδομένων}
\label{sec:step9-data}
% ------------------------------------------------------------------

Μετά την ολοκλήρωση του στησίματος, φορτώνονται τα δεδομένα του
\textlatin{Plegma Dataset}~\cite{plegma2024} στον \textlatin{Orion-LD}
μέσω του \textlatin{script} \textlatin{orion\_ld\_loading\_script.py}:

\selectlanguage{english}
\begin{lstlisting}[language=bash, caption={Data loading.}, stringstyle=\color{black}, keywordstyle=\color{black}]
cd Orion-LD
python orion_ld_loading_script.py
\end{lstlisting}
\selectlanguage{greek}

Το \textlatin{script} διαβάζει τα αρχεία \textlatin{CSV} του
\textlatin{dataset}, μετατρέπει κάθε γραμμή σε οντότητα
\textlatin{NGSI-LD} και τη δημοσιεύει στον \textlatin{Orion-LD}.
Δημιουργεί πρώτα τις στατικές οντότητες (\textlatin{Building},
\textlatin{Person}, \textlatin{Device}) και στη συνέχεια τις
χρονοσειρές (\textlatin{HouseholdElectricMeasurement},
\textlatin{EnvironmentalMeasurement}).

% ------------------------------------------------------------------
\section{Σύνοψη Διαδικασίας Στησίματος}
\label{sec:setup-summary}
% ------------------------------------------------------------------

Η διαδικασία στησίματος του \textlatin{Plegma Data Space} συνοψίζεται
στα ακόλουθα εννέα βήματα:

\begin{enumerate}
    \item Δημιουργία κρυπτογραφικών πιστοποιητικών και αντιγραφή
    στους κατάλληλους φακέλους.
    \item Παραμετροποίηση \textlatin{iSHARE Satellite} με τα
    αποτυπώματα πιστοποιητικών.
    \item Παραμετροποίηση \textlatin{Kong} με τα κρυπτογραφικά
    στοιχεία και τα \textlatin{endpoints}.
    \item Εκκίνηση \textlatin{containers} με τη σειρά
    \textlatin{Orion-LD} $\rightarrow$ \textlatin{Keyrock}
    $\rightarrow$ \textlatin{Kong}.
    \item Δημιουργία εφαρμογής \textlatin{OAuth} στον
    \textlatin{Keyrock}.
    \item Εφαρμογή \textlatin{patches} στον \textlatin{Keyrock}
    και επανεκκίνηση.
    \item Αρχική εγγραφή καταναλωτή για δημιουργία εσωτερικού
    \textlatin{UUID}.
    \item Δημιουργία πολιτικής εκχώρησης εξουσιοδότησης.
    \item Φόρτωση δεδομένων \textlatin{Plegma Dataset}.
\end{enumerate}

Στο επόμενο κεφάλαιο παρουσιάζονται αναλυτικά οι
\textlatin{end-to-end} ροές πρόσβασης στα δεδομένα, τόσο για
τον πάροχο όσο και για τον καταναλωτή, με βάση τα
\textlatin{sequence diagrams} που αντικατοπτρίζουν την πραγματική
συμπεριφορά του συστήματος.
